% Options for packages loaded elsewhere
\PassOptionsToPackage{unicode}{hyperref}
\PassOptionsToPackage{hyphens}{url}
%
\documentclass[
  11pt,
]{article}
\usepackage{amsmath,amssymb}
\usepackage{iftex}
\ifPDFTeX
  \usepackage[T1]{fontenc}
  \usepackage[utf8]{inputenc}
  \usepackage{textcomp} % provide euro and other symbols
\else % if luatex or xetex
  \usepackage{unicode-math} % this also loads fontspec
  \defaultfontfeatures{Scale=MatchLowercase}
  \defaultfontfeatures[\rmfamily]{Ligatures=TeX,Scale=1}
\fi
\usepackage{lmodern}
\ifPDFTeX\else
  % xetex/luatex font selection
\fi
% Use upquote if available, for straight quotes in verbatim environments
\IfFileExists{upquote.sty}{\usepackage{upquote}}{}
\IfFileExists{microtype.sty}{% use microtype if available
  \usepackage[]{microtype}
  \UseMicrotypeSet[protrusion]{basicmath} % disable protrusion for tt fonts
}{}
\makeatletter
\@ifundefined{KOMAClassName}{% if non-KOMA class
  \IfFileExists{parskip.sty}{%
    \usepackage{parskip}
  }{% else
    \setlength{\parindent}{0pt}
    \setlength{\parskip}{6pt plus 2pt minus 1pt}}
}{% if KOMA class
  \KOMAoptions{parskip=half}}
\makeatother
\usepackage{xcolor}
\usepackage[margin=1in]{geometry}
\usepackage{longtable,booktabs,array}
\usepackage{calc} % for calculating minipage widths
% Correct order of tables after \paragraph or \subparagraph
\usepackage{etoolbox}
\makeatletter
\patchcmd\longtable{\par}{\if@noskipsec\mbox{}\fi\par}{}{}
\makeatother
% Allow footnotes in longtable head/foot
\IfFileExists{footnotehyper.sty}{\usepackage{footnotehyper}}{\usepackage{footnote}}
\makesavenoteenv{longtable}
\setlength{\emergencystretch}{3em} % prevent overfull lines
\providecommand{\tightlist}{%
  \setlength{\itemsep}{0pt}\setlength{\parskip}{0pt}}
\setcounter{secnumdepth}{-\maxdimen} % remove section numbering
\usepackage{amsmath,amssymb,booktabs}
\ifLuaTeX
  \usepackage{selnolig}  % disable illegal ligatures
\fi
\IfFileExists{bookmark.sty}{\usepackage{bookmark}}{\usepackage{hyperref}}
\IfFileExists{xurl.sty}{\usepackage{xurl}}{} % add URL line breaks if available
\urlstyle{same}
\hypersetup{
  pdftitle={Evidence Against Authority: Costly Verification and the Audit of Unchanged Decisions},
  pdfauthor={Dr.~Jan-Hendrik Weinert},
  hidelinks,
  pdfcreator={LaTeX via pandoc}}

\title{Evidence Against Authority: Costly Verification and the Audit of
Unchanged Decisions}
\author{Dr. Jan-Hendrik Weinert}
\date{19 October 2026}


% Publication formatting edition, 19 October 2026. Scientific equations unchanged.
\usepackage{xurl,needspace,etoolbox}
\geometry{a4paper,margin=25.4mm}
\hypersetup{hidelinks,pdfauthor={Dr. Jan-Hendrik Weinert}}
\setlength{\parindent}{1.2em}
\setlength{\parskip}{3pt plus 1pt}
\setlength{\emergencystretch}{3em}
\hfuzz=0.2pt
\raggedbottom
\pagestyle{plain}
\makeatletter
\renewcommand{\@seccntformat}[1]{\csname the#1\endcsname\ifstrequal{#1}{section}{.}{}\quad}
\renewcommand{\maketitle}{%
  \begin{center}
  {\LARGE\bfseries\@title\par}
  \vspace{0.8em}
  {\normalsize\@author\par}
  \vspace{0.35em}
  {\small\@date\par}
  \end{center}
  \vspace{0.4em}}
\renewenvironment{abstract}{%
  \begin{center}\small\bfseries Abstract\end{center}
  \begin{list}{}{\setlength{\leftmargin}{0.5in}\setlength{\rightmargin}{0.5in}}
  \item[]\small\setlength{\parindent}{0pt}\ignorespaces
}{\end{list}\vspace{0.6em}}
\makeatother

\begin{document}
\maketitle
\begin{abstract}
Reputational resistance can reduce investigation even when managers use acquired evidence efficiently. We solve a Bayesian acquisition-and-correction game with an authenticated initial diagnosis and costly process audits. Career concerns depress private acquisition below its efficient level. Auditing unchanged decisions encourages investigation, whereas auditing reversals is redundant in equilibrium. For uniform investigation costs, we characterize the welfare-optimal targeted audit, including an explicit cost threshold and conditions for a unique partial audit. A private-ability counterexample shows why documentation need not restore efficient acquisition when information purchases also reveal type. The results apply to the stated authenticated-record and committed-audit institution; numerical examples are synthetic.
\end{abstract}
\section{1. Introduction}\label{introduction}

An unchanged decision is ambiguous. It can reflect confirming evidence,
justified confidence in an original diagnosis, a suppressed warning, or
a choice not to investigate. Organizational records often document the
final decision more readily than the search that preceded it.
Consequently, a policy aimed only at encouraging managers to reverse bad
decisions may miss the principal response to career concerns: fewer
costly checks are commissioned in the first place.

This article studies a manager whose earlier diagnostic accuracy affects
a Bayesian evaluator's assessment. A new signal can improve the final
action but can also contradict the original diagnosis. The manager
privately chooses whether to acquire that signal and then chooses either
final action. An auditor can retrieve a sealed record of acquisition
and, if acquired, the signal. The audit is costly and its rule is
announced before investigation. The evaluator observes the original
diagnostic output and the final action, with an audit supplying
additional process information before reputation is assessed.

The equilibrium specializes the established instrumental role of costly
information to an authenticated prior diagnosis. In a fixed-information
benchmark, a manager may retain an initially preferred action despite
understanding contradictory evidence. With a privately purchased, purely
instrumental signal and no ability information beyond the signal
history, such behavior cannot coexist with positive-cost acquisition in
equilibrium. If the continuation gain from acting on contradiction were
zero or negative, buying the signal would have no positive private
return. Managers who pay to investigate must therefore strictly prefer
to use its corrective content. Career concerns remain inefficient, but
they reduce investigation rather than the conditional correction
probability.

The second result concerns audit design. Once all purchased
contradictory evidence is used, a reversal already identifies the
acquisition of a contradictory signal. An audit of that reversal
supplies no new information about initial diagnostic ability. An
unchanged decision, by contrast, pools an uninvestigated diagnosis with
an investigated and confirmed one. Auditing this pooled event makes
investigation more attractive. The model determines acquisition and the
audit workload jointly: a successful policy generates more reversals,
reducing the fraction of decisions that need to be inspected. We solve
the optimal audit within a clearly specified feasible class and obtain
an explicit cost threshold for whether any audit is worthwhile.

The economic trade-off is neither an assumed disincentive to effort
under transparency nor an assumed preference for partial disclosure.
Disclosure increases acquisition in the main environment. Partial audit
arises because information and audit resources have opportunity costs.
Under a stated regularity condition, any strictly positive audit cost
makes complete auditing of unchanged decisions excessive; a sufficiently
large cost makes no auditing optimal. Outside that condition, a finite
candidate characterization still gives the global optimum without
presuming an interior solution.

The results are deliberately conditional. The initial diagnostic output
is authenticated, though its endorsement can be chosen strategically.
The manager initially lacks private information about own diagnostic
ability. We give a fully solved counterexample with known ability in
which acquisition itself becomes a bad signal and remains inefficient
under full documentation. That counterexample identifies a boundary of
the main theorem rather than adding an unsolved general model. A
separate protocol proposes a compact test of evidence acquisition with
signal quality held fixed. The article establishes theoretical results,
not empirical effects on managerial behavior or firm profits.

\section{2. Related literature and the specific
contribution}\label{related-literature-and-the-specific-contribution}

Prendergast and Stole (1996) show that career concerns distort a
sequence of managerial investment choices, including later conservatism.
Their result establishes the reputational mechanism underlying our
fixed-information benchmark. Our equilibrium adds an authenticated
diagnostic record and a unique endogenous acquisition cutoff in the
stated parameter region. The positive-cost instrumental-acquisition step
has a close antecedent discussed below and is not offered as
independently original. Initial diagnostic ability is not an ability to
receive better corrective signals; the later signal has common
precision.

Prat (2005) distinguishes disclosure of an action from disclosure of its
consequence. In his May 2004 working paper, Proposition 4 separates
effects on decision discipline from effects on sorting; Proposition 5
shows that greater consequence observability permits more action
observability while preserving a separating equilibrium. That
proposition restricts attention to pooling and separating equilibria and
does not establish a general sorting comparison. Our audit retrieves
acquisition and its signal, not a realized consequence. Actions remain
public, and the planner here assigns no direct value to sorting. The new
audit theorem solves a resource-cost problem over action-contingent
trail inspections after acquisition becomes endogenous; it does not
claim that complementarity between accountability records was previously
unknown.

Levy (2007) studies the interaction of voting rules, transparency, and
reputation in committees. Concealing individual actions need not
eliminate reputational distortions. Our result concerns a single manager
and the allocation of inspection across corrected and unchanged
decisions. It makes no claim about optimal voting or the quality of
collective information aggregation.

Bar-Isaac (October 2011 manuscript, Sections 2-4) studies one safe and
two risky actions, costly binary effort, observed outcomes, and an
action-linked premium \(R\). The opening of his Proposition 1 proof
excludes costly effort followed by the information-insensitive safe
action. His Proposition 1 compares existence of effort equilibria across
action-observation regimes according to the sign of \(R\); Section 4
also discusses direct observation of information. Thus the
instrumental-acquisition logic used in our equation (6) has an explicit
antecedent. We do not claim it as original. The cited manuscript is
distinguished bibliographically from the 2012 final publication.

The specific additional content here concerns an authenticated diagnosis
and a feasible audit technology. Retention here yields
\(1\{s=\theta\}\), not an intrinsically constant payoff. Contradiction
informs the evaluator about \emph{earlier} diagnostic ability, while the
corrective signal's precision is independent of that ability. Under the
independent-cost, unknown-own-type information structure of Sections
3.1-3.2 and the bounds (D) in Section 3.3, Theorem 1 derives a unique
acquisition allocation without selecting label symmetry. Equations
(4)-(5) determine both the cutoff and the endogenous posterior wedge;
equation (6) is an environment-specific application of the antecedent
instrumental logic. The theorem is not a general claim that paying for
information always ensures efficient use.

The audit result is the main incremental design content. Actions remain
public and final outcomes are withheld until after career assessment.
Within the committed two-probability signal-trail technology of Section
3.2, Proposition 1 in Section 5.1 proves that reversal inspection is
redundant, while retention inspection distinguishes no acquisition from
confirmation. Its endogenous workload is \(\tau_N(1-dq)\) in equation
(8). For uniform costs, equations (10)-(12) solve the global feasible
optimum; Theorem 3 and equations (13)-(15) give the additional condition
and explicit cost threshold for unique partial auditing versus no
auditing. These are not consequences of equation (6) alone: they require
the diagnostic posterior identities, the specified inspection contents,
real inspection costs, and the feasible policy restrictions. Neither an
unrestricted audit optimum nor a direct sorting benefit is claimed.

Gersbach and Hahn (2012) study costly information and transparency in
committees, including implications for effort and subsequent
composition. Together with the antecedent above, this rules out a broad
originality claim based on endogenizing information. Our bounded
contribution is the jointly solved diagnostic-record equilibrium,
retention-trail targeting, and endogenous-workload audit-cost
characterization. Prat's consequence records remain distinct from the
signal trails inspected here.

Ishida and Suen (2026) show that quitting can signal an ability to
recognize project quality; informed agents respond efficiently to news
while other agents distort their quitting times. Their information
structure cautions against treating revision as an intrinsically
negative ability signal. In our main model, contradiction predicts lower
\emph{initial} diagnostic accuracy, whereas ability does not affect the
precision of corrective evidence. The private-ability counterexample
then shows how endogenous selection into acquisition can revive an
adverse reputational signal even with complete records.

Hoffman, Kahn, and Li (2018) motivate distinguishing informative
managerial discretion from harmful departures from testing
recommendations. Dietvorst, Simmons, and Massey (2015) and Yaniv and
Kleinberger (2000) identify distinct problems of algorithm reliance and
advice discounting. The mechanism here does not require mistaken private
beliefs, distrust of a signal, or an algorithm. Its identifying
prediction is about whether evidence is acquired, conditional on an
unchanged and known information technology.

The contribution is thus a solved conjunction of hidden costly
corrective search, endogenous correction, and resource-constrained audit
targeting. The claim is limited to the propositions proved here, rather
than priority over every model of strategic ignorance or transparency.

\section{3. Environment}\label{environment}

\subsection{3.1 Signals, private cost, and
actions}\label{signals-private-cost-and-actions}

The binary state \(\theta\in\{0,1\}\) has prior one half. Diagnostic
ability \(T\in\{H,L\}\) is independent of \(\theta\), with
\(P(T=H)=\mu\in(0,1)\). An authenticated diagnostic output
\(s\in\{0,1\}\) satisfies \(P(s=\theta\mid T=t)=p_t\), with
\(1/2<p_L<p_H<1\). Both the evaluator and manager initially lack
information about \(T\) beyond the common prior. In particular,
observing the label \(s\) does not change that prior.

After seeing \(s\), the manager privately observes a cost \(c\),
independently drawn from a continuous distribution \(F\) on \([0,C]\),
with positive density on its interior. The manager chooses acquisition
\(e\in\{0,1\}\). Acquisition costs \(c\) and produces \(z\in\{0,1\}\)
with accuracy \(r\), independently of \((s,T)\) conditional on
\(\theta\). If \(e=0\), no \(z\) is available to the manager. If
\(e=1\), its actual realization is irreversibly recorded. There is no
forgery or record deletion.

The final action is unrestricted: \(a\in\{0,1\}\), whether or not
evidence was acquired and whether or not it confirms \(s\). The state is
not revealed to the evaluator before reputation is paid. We call \(a=s\)
retention and \(a\ne s\) reversal; these terms refer to the diagnostic
output, not to an arbitrary public message.

The manager's expected utility is

\[\mathbb E[1\{a=\theta\}+\lambda m-ec],\tag{1}\]

where \(m\) is the evaluator's posterior probability that \(T=H\) and
\(\lambda>0\). The evaluator follows Bayes' rule and has no independent
strategic wage-setting decision. Career payoff is a reduced form. The
planner values correct action, incurs audit cost, and internalizes real
investigation cost when evaluating total surplus. We distinguish this
welfare objective from a firm's accounting profit that excludes the
manager's effort cost.

\subsection{3.2 Committed audit
technology}\label{committed-audit-technology}

The planner chooses \((\tau_N,\tau_R)\in[0,1]^2\), before acquisition,
where \(\tau_N\) is the probability of auditing a retention and
\(\tau_R\) the probability of auditing a reversal. The rule is public.
The random audit realization occurs after action and is conditionally
independent of all private information given retention or reversal. The
evaluator always observes \((s,a)\) and the audit flag. An audit
retrieves either the record \(e=0\), or the record \((e=1,z)\). Cost is
\(\kappa>0\) per actual inspection, even when no evidence was purchased.
Audit information arrives before reputation is assessed.

The planner cannot directly purchase the manager's evidence, impose
acquisition mandates or fines, pay type-contingent rewards, observe
\(c\), or choose arbitrary garblings. These restrictions define the
feasible audit problem. A different institutional instrument can
dominate this class; our optimum is not an unrestricted mechanism-design
optimum. Audit commitment is essential, since a later evaluator can
prefer more information than the ex ante policy permits.

We use Bayesian sequentially consistent assessments: strategies and
their perturbations are measurable with respect to the manager's actual
information. Because the manager does not observe \(T\), trembles cannot
condition directly on it. At a fully documented history, any compatible
perturbation has the same distribution of private cost conditional on
\(T\). Thus an off-path action cannot arbitrarily change ability beliefs
conditional on a complete record. This is a consequence of the
information structure, proved in Appendix A. It differs from assigning
an off-path reputation to an otherwise unverifiable reversal.

\subsection{3.3 Derived quantities and the parameter
domain}\label{derived-quantities-and-the-parameter-domain}

Let

\[p=\mu p_H+(1-\mu)p_L,\quad h=r-p>0,\]
\[d_t=p_t(1-r)+(1-p_t)r,\quad d=\mu d_H+(1-\mu)d_L,\]
\[m_C=\frac{\mu d_H}{d},\quad m_F=\frac{\mu(1-d_H)}{1-d},\quad A=\mu-m_C>0,\]
\[g=\frac{h}{d},\quad B=m_F-m_C=\frac{A}{1-d},\quad L=\lambda d A.\tag{2}\]

Here \(d\) is the probability of contradiction conditional on
acquisition. With both signals, the accuracy advantage of following
\(z\) at contradiction is \(g\). A confirming pair favors retention more
strongly than \(s\) alone. Efficient acquisition occurs at \(c<h\), so
the efficient acquisition rate is \(q^{FB}=F(h)\).

The main equilibrium result uses the open parameter region

\[C>h,\qquad \lambda B<2p-1,\qquad L<h.\tag{D}\]

Every structurally compatible observation has a type posterior in
\([m_C,m_F]\), as shown in Appendix A. The first reputational bound
makes an uninformed reversal strictly worse even under the greatest
compatible reputation improvement. It also rules out reversing
confirming evidence. The second ensures that sufficiently cheap
investigation is profitable even if an off-path reversal received the
lowest compatible reputation \(m_C\). These are sufficient, economically
interpretable restrictions, not claims about all career-concern
environments. Together they rule out complete acquisition as well as
zero acquisition. No label-symmetry selection is imposed.

\section{4. Equilibrium: reputation moves the distortion into
acquisition}\label{equilibrium-reputation-moves-the-distortion-into-acquisition}

\subsection{4.1 A fixed-information
benchmark}\label{a-fixed-information-benchmark}

If everyone receives \(z\) at no choice or cost, a public reversal can
reveal a contradictory event associated with lower initial diagnostic
ability. With neither audit nor outcomes, a candidate strategy that
reverses a fraction \(\rho\) of contradictions gives

\[m_N(\rho)=\frac{\mu(1-\rho d_H)}{1-\rho d},\quad
m_N(\rho)-m_C=\frac{A}{1-\rho d}.\tag{3}\]

Within evidence-following strategies, incomplete correction can be
supported when its accuracy gain equals the reputational cost. This is a
standard reputational-conservatism benchmark. Neither its correction
formula nor an arbitrary belief at a zero-revision boundary is a main
result of this article. The endogenous acquisition game instead has both
final actions on path under (D), eliminating that boundary problem.

\subsection{4.2 Authenticated-diagnosis
equilibrium}\label{authenticated-diagnosis-equilibrium}

\textbf{Theorem 1 (unique costly-search equilibrium).} Assume (D) and
the information structure of Section 3. For every committed audit rule,
all Bayesian sequentially consistent equilibria have the same outcome,
except for a probability-zero acquisition-cost tie:

\begin{enumerate}
\def\labelenumi{\arabic{enumi}.}
\tightlist
\item
  An uninformed manager retains \(s\). An informed manager retains after
  confirmation and reverses after contradiction.
\item
  The manager acquires precisely when \(c<c^*\). The acquisition rate
  \(q=F(c^*)\) is the unique solution to
\end{enumerate}

\[c^*=h-\frac{(1-\tau_N)L}{1-dq},\qquad q=F(c^*).\tag{4}\]

\begin{enumerate}
\def\labelenumi{\arabic{enumi}.}
\setcounter{enumi}{2}
\tightlist
\item
  Both retention and reversal occur with positive probability.
  Conditional on either diagnostic label,
\end{enumerate}

\[m_R=m_C,\quad m_N(q)=\frac{\mu(1-qd_H)}{1-qd},\quad D(q)=\frac{A}{1-qd}.\tag{5}\]

\begin{enumerate}
\def\labelenumi{\arabic{enumi}.}
\setcounter{enumi}{3}
\tightlist
\item
  The rate \(q\) is independent of \(\tau_R\), increases strictly in
  \(\tau_N\) below one, satisfies \(0<q<q^{FB}\) for \(\tau_N<1\), and
  equals \(q^{FB}\) at \(\tau_N=1\).
\end{enumerate}

\emph{Economic argument.} The instrumental step below is an application
of the antecedent identified in Section 2, not a separate priority
claim. If buying information were followed by mixing over correction at
an exact continuation indifference, its incremental private value before
the cost would be zero. A positive acquisition cost therefore rules out
that combination. Managers can be deterred from investigating, but those
who pay to investigate must prefer to implement its useful
contradiction. All retention histories then pool absence of
investigation with confirming evidence. Equation (4) determines how the
expected reputation loss of leaving that pool reduces willingness to
investigate.

\emph{Proof.} Every observation coarsens the three primitive information
events: no acquisition, acquired confirmation, and acquired
contradiction. Their type posteriors are \(\mu,m_F,m_C\), so consistency
bounds every posterior by \([m_C,m_F]\). An uninformed reversal loses
\(2p-1\) in expected accuracy, more than the largest compatible
reputation gain \(\lambda B\). A reversal after confirmation loses still
more. These actions are strictly dominated, regardless of off-path
beliefs. At a fully documented history, the action supplies no further
type information. An unaudited reversal, if on path, must therefore
originate in acquired contradictory evidence.

First rule out \(q=0\). Retention then yields posterior \(\mu\). A
manager can buy evidence and reverse exactly at contradiction. Even
giving an unaudited reversal posterior \(m_C\), the gain from this
deviation is at least \(h-\lambda d(\mu-m_C)-c=h-L-c>0\) for
sufficiently small \(c\), by (D). Acquired evidence has positive
probability in any equilibrium.

The correction incentive is common to all cost types once the cost is
sunk. If correction occurred with probability \(\rho\), the net gain
from acquisition would be

\[d\rho\{g-(1-\tau_N)\lambda D\}-c.\tag{6}\]

If the bracket is negative, correction never occurs. If zero,
acquisition still loses \(c>0\). Both contradict positive acquisition.
Hence the bracket is strictly positive and every acquired contradiction
is corrected. This determines the on-path posteriors (5).

An audited reversal also yields \(m_C\), so \(\tau_R\) has no effect. An
audited retention gives posterior \(\mu\) when \(e=0\) and \(m_F\) after
confirmation. The identity \(dm_C+(1-d)m_F=\mu\) implies that acquiring
and then following the signal changes expected reputation, relative to
not acquiring, by \(-(1-\tau_N)dD(q)\). Adding accuracy gain \(h\) and
subtracting cost gives (4).

The right side of (4) decreases with \(q\), whereas \(F^{-1}(q)\)
increases. At \(q=0\) it is \(h-(1-\tau_N)L>0\); at \(q=F(h)\) it is at
most \(h\), strictly below for \(\tau_N<1\). A unique solution follows
by continuity. At that solution the correction gain is \(c^*/d>0\), so
all claimed actions and the acquisition cutoff are sequential best
responses. These arguments apply separately for each \(s\) and produce
the same unique solution, rather than selecting a symmetric equilibrium.
Complete documentation pins down off-path action beliefs as explained in
Appendix A. \(\square\)

\textbf{Corollary 1 (conditional correction can conceal an
organizational distortion).} Within this environment, the conditional
probability of using acquired contradictory evidence is one under every
audit rule. Nevertheless, with \(\tau_N<1\), real investment in evidence
is inefficiently low. A comparison restricted to informed decision
makers misses the distortion.

This is a selection result, not a prescription to measure revision rates
unconditionally. Total reversals are \(qd\) and expected accuracy is
\(p+qh\). They can both increase when auditing encourages investigation,
even though conditional evidence use never changes.

\subsection{4.3 Implicit acquisition comparative
statics}\label{implicit-acquisition-comparative-statics}

For a positive density \(f\) at \(c^*\), differentiating (4) gives

\[\frac{dq}{d\tau_N}=\frac{L/(1-dq)}{1/f(c^*)+(1-\tau_N)Ld/(1-dq)^2}>0.\tag{7}\]

A lower direct investigation cost can similarly raise acquisition, but
it also raises the reputation of unchanged decisions by making their
confirmation interpretation more credible. The denominator captures this
equilibrium feedback. Greater acquisition does not leave the reputation
wedge fixed. This is why replacing \(\lambda\) with a
disclosure-adjusted constant in a fixed-information correction formula
does not solve the present model.

\section{5. Welfare-optimal targeted
documentation}\label{welfare-optimal-targeted-documentation}

\subsection{5.1 Why reversals need no inspection in this
class}\label{why-reversals-need-no-inspection-in-this-class}

\textbf{Proposition 1 (audit targeting).} For \(\kappa>0\), every
welfare-optimal policy in the class of Section 3 has \(\tau_R=0\).

\emph{Proof.} Theorem 1 shows that acquired contradictions are always
corrected and that reversal already implies acquisition and
contradiction. Inspection of a reversal therefore leaves the ability
posterior at \(m_C\). Changing \(\tau_R\) leaves every acquisition and
action incentive unchanged. Since \(qd>0\), a positive \(\tau_R\) uses
strictly positive audit resources with no information or decision
benefit in this model. \(\square\)

Targeting is not an assumption that unchanged decisions are suspicious
in an everyday sense. It follows from the informativeness of equilibrium
reversal and from the audit's limited contents. An audit that can
discover fraud, undocumented private signals, or manipulation may make
reversal inspections valuable; those benefits are absent here.

Write \(\tau=\tau_N\). Welfare relative to the original unaudited
diagnosis is

\[V(\tau)=h q-\int_0^{c^*}c\,dF(c)-\kappa\tau(1-dq),\tag{8}\]

where \(q\) and \(c^*\) are determined by (4). The audit rate is
\(\tau(1-dq)\), not \(\tau\): the policy does not inspect reversals.
Gross firm decision value is \(p+qh\). Total surplus also deducts real
investigation costs. Ex ante mean reputation is \(\lambda\mu\) under
every Bayesian policy and therefore drops out of the comparison. Career
assessments redistribute rents across histories; they are not themselves
a net resource gain in this linear-payoff environment.

\subsection{5.2 Exact global solution under uniform investigation
costs}\label{exact-global-solution-under-uniform-investigation-costs}

Suppose \(F(c)=c/C\) on \([0,C]\). Define \(q_0=q(0)\) and \(q_f=h/C\).
They satisfy

\[Cq_0=h-\frac{L}{1-dq_0},\quad 0<q_0<q_f<1.\tag{9}\]

The audit rate required to implement acquisition probability
\(q\in[q_0,q_f]\) is

\[\tau(q)=1-\frac{(h-Cq)(1-dq)}{L}.\tag{10}\]

It is strictly increasing and maps the feasible interval onto \([0,1]\).
Let

\[H(q)=\tau(q)(1-dq),\quad
V(q)=h q-\frac{Cq^2}{2}-\kappa H(q).\tag{11}\]

\textbf{Theorem 2 (global finite characterization).} The welfare-optimal
feasible audit is obtained by evaluating (11) at \(q_0\), \(q_f\), and
the real roots in \([q_0,q_f]\) of

\[0=h+\kappa d-Cq-\frac{\kappa}{L}
\{C+2dh-(4Cd+2d^2h)q+3Cd^2q^2\}.\tag{12}\]

The maximizing candidate, mapped through (10), is globally optimal. Ties
are genuine policy multiplicity and can be resolved by the smaller audit
rate. No local first-order condition alone is used to claim optimality.

\emph{Proof.} Equation (10) implements each and only each feasible
acquisition rate by Theorem 1. Substitution makes (11) a cubic
polynomial on a compact interval. Its global maximum exists and lies at
an endpoint or a zero of its derivative, which is (12). Proposition 1
has already removed the redundant audit instrument. \(\square\)

The characterization is valid without presuming that welfare is concave,
audit demand is increasing, or the optimum is interior. A sharper
economically interpretable result follows in a nonempty parameter
region.

\subsection{5.3 An explicit audit-cost
threshold}\label{an-explicit-audit-cost-threshold}

Let \(t_0=1-dq_0\) and \(x_0=Ld/(Ct_0^2)\). Define

\[x_c=\frac{\sqrt{17}-3}{4},\qquad
\kappa_c=\frac{L^2}{t_0(Ld+Ct_0^2)}.\tag{13}\]

\textbf{Theorem 3 (unique partial audit or no audit).} In addition to
(D), suppose \(x_0<x_c\). Then the optimal policy is unique:

\begin{itemize}
\tightlist
\item
  At \(\kappa=0\), \(q^*=q_f\) and \(\tau^*=1\), up to the costless
  redundant reversal-audit instrument.
\item
  For \(0<\kappa<\kappa_c\), \(q_0<q^*<q_f\) and \(0<\tau^*<1\). The
  larger root of (12) in the feasible interval is the unique optimum.
\item
  For \(\kappa\geq\kappa_c\), \(q^*=q_0\) and \(\tau^*=0\).
\end{itemize}

The optimal acquisition and audit probabilities decrease strictly with
\(\kappa\) in the interior region. Complete auditing of unchanged
decisions is inefficient at every strictly positive audit cost in this
parameter region.

\emph{Proof.} With \(t=1-dq\),

\[H'=-d+\frac{Ct^2+2d(h-Cq)t}{L},\]
\[H''=-\frac{4Cd t+2d^2(h-Cq)}{L}<0,\quad H'''=\frac{6Cd^2}{L}>0.\tag{14}\]

At the no-audit endpoint, (9) implies

\[H'(q_0)=d+Ct_0^2/L,\quad V'(q_0)=L/t_0-\kappa H'(q_0).\]

Its zero is exactly \(\kappa_c\). Since \(1-dq_f=t_0(1-x_0)\),

\[H'(q_f)=d\{(1-x_0)^2/x_0-1\}>0\]

under \(x_0<x_c<(3-\sqrt5)/2\). As \(H''<0\), \(H'>0\) over the entire
feasible interval.

At \(\kappa=\kappa_c\),

\[\frac{\kappa_c[-H''(q_0)]}{C}=\frac{4x_0+2x_0^2}{1+x_0}<1.\tag{15}\]

Because \(H''\) increases with \(q\), this makes \(V''=-C-\kappa H''<0\)
throughout the interval for all \(0\leq\kappa\leq\kappa_c\). When
\(0<\kappa<\kappa_c\), \(V'(q_0)>0\) while
\(V'(q_f)=-\kappa H'(q_f)<0\), giving a unique interior maximum. At
\(\kappa_c\), strict concavity and \(V'(q_0)=0\) give unique optimum
\(q_0\). For larger \(\kappa\), relative welfare at any \(q>q_0\) falls
by the additional positive amount \((\kappa-\kappa_c)[H(q)-H(q_0)]\), so
no auditing remains uniquely optimal. Finally, implicit differentiation
of the interior first-order condition gives \(dq^*/d\kappa=H'/V''<0\),
and (10) is increasing. \(\square\)

Partial auditing therefore follows from a derived trade-off between
additional socially valuable investigation and real inspection cost. The
regularity condition concerns the geometry of the audit workload, not an
imposed sign of effort's response. Theorem 2 supplies the global answer
when that condition fails.

\subsection{5.4 What the optimum does not
cover}\label{what-the-optimum-does-not-cover}

The planner does not receive Prat's convex benefit from identifying and
replacing high- or low-ability managers. Adding a sorting payoff can
make extra information valuable even when final accuracy does not
change, and it can change optimal audit targeting. Likewise, a firm that
excludes the manager's investigation cost from its objective maximizes
\(hq-\kappa H(q)\), not (11). Its optimum is obtained from the same
endpoint comparison with derivative \(h-\kappa H'(q)\), and need not
share Theorem 3's interior/no-audit threshold. The article's normative
result is explicitly about total surplus within the stated audit
technology.

\section{6. Information about own ability and strategic initial
endorsement}\label{information-about-own-ability-and-strategic-initial-endorsement}

\subsection{6.1 A solved private-ability
counterexample}\label{a-solved-private-ability-counterexample}

The main theorem does not extend merely by replacing a common accuracy
with type-specific accuracies. If a manager knows own ability,
acquisition can reveal that the original diagnosis was weak. Independent
documentation then cannot wash out the type information in the
acquisition decision itself.

\textbf{Proposition 2 (full trails need not restore efficient
acquisition).} Consider \(\mu=1/2\), \(p_H=.85\), \(p_L=.60\),
\(r=.80\), a common acquisition cost \(c=.05\), career weight
\(\lambda=.20\), and full documentation after every final action.
Suppose the manager knows \(T\). There is a sequentially rational
equilibrium in which high types never acquire, low types acquire with
probability \(2/3\), every acquired low type follows \(z\), and every
unacquired type retains \(s\). Efficient acquisition would have all low
types acquire and no high types acquire.

\emph{Proof.} In the proposed equilibrium, acquisition reveals \(L\)
with probability one, so its posterior is zero for every realized signal
and action on path. No acquisition has posterior

\[m_0=\frac{1/2}{1/2+(1/2)(1-2/3)}=3/4.\]

A low type gains \(.80-.60=.20\) in accuracy by acquiring and following
\(z\). Its utility difference is \(.20-.05-.20(3/4)=0\), so mixing is
optimal. After acquisition, both signal realizations have posterior zero
and following \(z\) strictly improves accuracy at contradiction. A high
type's information benchmark retains \(s\) even after contradiction,
because \(.85>.80\); acquiring therefore provides no accuracy gain,
costs \(.05\), and lowers its reputation from \(.75\) to zero. It never
acquires.

Uninformed low-type reversal loses \(.20\) in accuracy; even the highest
off-path posterior one can improve its reputation payoff by at most
\(.20(1-.75)=.05\). It therefore retains. For high types the accuracy
loss is \(.70\), so retention is also strict. At acquired histories
every signal/action pair used by a low type has positive probability.
Other action histories can be assigned posterior zero by consistent
perturbations in which high-type acquisition is arbitrarily rarer than
low-type action mistakes. None benefits either type. Thus all
acquisition and action deviations are checked. Without career concerns,
low-type net acquisition value \(.20-.05=.15\) is positive, while
high-type value is negative. \(\square\)

This finite example does not characterize all private-ability
equilibria. It proves that Theorem 1's complete-documentation conclusion
is not robust to private type information. In particular,
type-independent acquisition is the substantive reason why an audited
absence of acquisition receives prior reputation in the main
environment.

\subsection{6.2 Endogenous endorsement with an authenticated diagnostic
output}\label{endogenous-endorsement-with-an-authenticated-diagnostic-output}

\textbf{Proposition 3 (strategic endorsement conditional on a raw
record).} Add a freely chosen initial recommendation \(a_0\in\{0,1\}\)
after observing \(s\), before the cost draw, and a payoff
\(\beta1\{a_0=\theta\}\), with \(\beta>0\). Keep \(s\) authenticated and
visible to the evaluator, and condition audit policy on final action
relative to \(s\). Under (D), the equilibrium of Theorem 1 remains
unchanged, and \(a_0=s\) is the unique initial recommendation.

\emph{Proof.} The manager's only information at the endorsement stage is
\(s\), not own type or future cost. Conditional on \(s\), a
recommendation or its tremble has no further information about \(T\).
Subsequent signals, cost distribution, and audit technology do not
depend on the recommendation. The unique continuation of Theorem 1
consequently gives the same payoff after either \(a_0\). The initial
accuracy gain from choosing \(s\) is \(\beta(2p-1)>0\). \(\square\)

The raw-record qualification is essential. If the manager can suppress
or falsify \(s\), or if audit classification uses a manipulable initial
message instead of the diagnostic output, Proposition 3 does not apply.
We have endogenized endorsement, not solved arbitrary strategic
production of the first signal. The main new hidden action is the
privately chosen production of corrective evidence.

\section{7. Numerical verification and implications for
measurement}\label{numerical-verification-and-implications-for-measurement}

The code implements the analytic equilibrium, all feasible terminal
action plans, and the global audit solution. It uses exact expectations
over type, state, and corrective signals rather than treating synthetic
decisions as empirical observations. The calculations below are
illustrative design inputs, not a calibration.

Set \(\mu=.5\), \(p_H=.85\), \(p_L=.60\), \(r=.80\), \(\lambda=1.5\),
and \(C=.15\). These values satisfy (D). Then \(p=.725\), \(h=.075\),
\(d=.365\), \(L=.05625\), \(q_0=.109403\), and \(q_f=.5\). The audit
regularity measure is \(x_0=.148498<x_c\), and \(\kappa_c=.02075475\).

\begin{longtable}[]{@{}
  >{\raggedleft\arraybackslash}p{(\columnwidth - 8\tabcolsep) * \real{0.2000}}
  >{\raggedleft\arraybackslash}p{(\columnwidth - 8\tabcolsep) * \real{0.2000}}
  >{\raggedleft\arraybackslash}p{(\columnwidth - 8\tabcolsep) * \real{0.2000}}
  >{\raggedleft\arraybackslash}p{(\columnwidth - 8\tabcolsep) * \real{0.2000}}
  >{\raggedleft\arraybackslash}p{(\columnwidth - 8\tabcolsep) * \real{0.2000}}@{}}
\toprule\noalign{}
\begin{minipage}[b]{\linewidth}\raggedleft
Audit cost \(\kappa\)
\end{minipage} & \begin{minipage}[b]{\linewidth}\raggedleft
Optimal \(\tau_N\)
\end{minipage} & \begin{minipage}[b]{\linewidth}\raggedleft
Acquisition \(q\)
\end{minipage} & \begin{minipage}[b]{\linewidth}\raggedleft
Actual inspection rate
\end{minipage} & \begin{minipage}[b]{\linewidth}\raggedleft
Net welfare gain over no audit
\end{minipage} \\
\midrule\noalign{}
\endhead
\bottomrule\noalign{}
\endlastfoot
0.000 & 1.000000 & 0.500000 & 0.817500 & 0.011442 \\
0.005 & 0.881826 & 0.447044 & 0.737937 & 0.007542 \\
0.010 & 0.721369 & 0.378751 & 0.621644 & 0.004123 \\
0.020 & 0.088175 & 0.139694 & 0.083679 & 0.000032 \\
0.030 & 0.000000 & 0.109403 & 0.000000 & 0.000000 \\
\end{longtable}

These are normalized accuracy and resource-cost units, with no empirical
calibration. At \(\kappa=.01\), gross accuracy rises by approximately
.02020 relative to no audit, while investigation and inspection reduce
the net welfare gain to .004123. Every investigator corrects every
contradiction under both policies. For comparison, the fixed-information
benchmark with the same signal parameters and \(\lambda=1.5\) has
\(\rho=.684932\) within evidence-following strategies. The new
equilibrium has conditional correction one but acquisition only .109403
without audit. The contrast concerns the location of the distortion, not
a welfare ranking between environments with different information costs.

For 400 independently generated admissible parameter configurations, the
reproduction program checks acquisition indifference, evaluates 4,800
complete terminal-plan deviations at costs on both sides of the cutoff,
and compares the analytic policy against 1,001 independently evaluated
audit probabilities per configuration. It also checks the threshold
theorem and the equilibrium under two nonuniform cost distributions. The
maximum acquisition-equation residual is below \(9\times10^{-17}\). The
maximal analytic advantage over the finite audit grid is about
\(5.54\times10^{-9}\), consistent with grid discretization. These
computations supplement, and do not substitute for, the proofs.

The principal measurement implication is that an audit policy can change
the composition of informed decision makers. The unconditional causal
effect of an announced policy on acquisition is well defined under
random assignment. Conditioning on realized acquisition generally
destroys that randomization. A near-perfect conditional correction rate
can coexist with too little investigation; an unchanged final action is
not by itself a measure of evidence rejection.

A separate design supplement proposes a two-arm,
evidence-technology-held-fixed test of paid acquisition. It retains a
compulsory-evidence validation task to distinguish acquisition from
comprehension, without making that selected conditional comparison the
primary estimand. The earlier eight-cell factorial is an optional,
superseded extension rather than the operational primary study. No
experiment has been run.

\section{8. Conclusion}\label{conclusion}

The instrumental role of costly information is an antecedent, rather
than this article's independent originality claim. In the
authenticated-diagnosis environment solved here, private positive-cost
acquisition changes the equilibrium expression of reputational
resistance: managers who investigate use contradictory evidence
efficiently, but too few investigate. The welfare loss is consequently
upstream of the observed correction decision.

This result leads to a precise audit implication. Within a committed
action-contingent trail-inspection technology, reversals need no
inspection because they already reveal the informative event. Inspecting
unchanged decisions increases investigation, and its endogenous workload
can be incorporated into an exact welfare solution. A sufficient
regularity condition delivers a unique partial audit below a stated cost
threshold and no audit above it.

These conclusions require authenticated first diagnoses, the absence of
private own-ability information in the main game, and the specified
career and welfare objectives. The private-ability counterexample shows
why documentation alone is not a universal remedy. The completed
theoretical results identify a particular mechanism and a feasible
policy trade-off; whether organizations exhibit that mechanism remains
an empirical question.


{
\small
\setlength{\parindent}{0pt}
\section*{References}\label{references}

Bar-Isaac, H. 2011. ``Transparency, Career Concerns, and Incentives for
Acquiring Expertise.'' Author manuscript, October.
\href{https://www.researchgate.net/publication/227378148_Transparency_Career_Concerns_and_Incentives_for_Acquiring_Expertise}{Author-uploaded
manuscript}.

Bar-Isaac, H. 2012. ``Transparency, Career Concerns, and Incentives for
Acquiring Expertise.'' \emph{The B.E. Journal of Theoretical Economics}
12(1), Article 4. \url{https://doi.org/10.1515/1935-1704.1796}.

Dietvorst, B. J., J. P. Simmons, and C. Massey. 2015. ``Algorithm
Aversion: People Erroneously Avoid Algorithms After Seeing Them Err.''
\emph{Journal of Experimental Psychology: General} 144(1): 114-126.
\url{https://doi.org/10.1037/xge0000033}.

Gersbach, H., and V. Hahn. 2012. ``Information Acquisition and
Transparency in Committees.'' \emph{International Journal of Game
Theory} 41(2): 427-453. \url{https://doi.org/10.1007/s00182-011-0295-5}.

Hoffman, M., L. B. Kahn, and D. Li. 2018. ``Discretion in Hiring.''
\emph{Quarterly Journal of Economics} 133(2): 765-800.
\url{https://doi.org/10.1093/qje/qjx042}.

Ishida, J., and W. Suen. 2026. ``Signaling Vision: Knowing When to
Quit.'' \emph{International Economic Review} 67(3): 997-1013.
\url{https://doi.org/10.1111/iere.70066}.

Levy, G. 2007. ``Decision Making in Committees: Transparency,
Reputation, and Voting Rules.'' \emph{American Economic Review} 97(1):
150-168. \url{https://doi.org/10.1257/aer.97.1.150}.

Prat, A. 2005. ``The Wrong Kind of Transparency.'' \emph{American
Economic Review} 95(3): 862-877.
\url{https://doi.org/10.1257/0002828054201297}. The corresponding May 26, 2004
working paper is available at
\url{https://eprints.lse.ac.uk/24712/1/dp498.pdf}.

Prendergast, C., and L. Stole. 1996. ``Impetuous Youngsters and Jaded
Old-Timers: Acquiring a Reputation for Learning.'' \emph{Journal of
Political Economy} 104(6): 1105-1134. \url{https://doi.org/10.1086/262055}.

Yaniv, I., and E. Kleinberger. 2000. ``Advice Taking in Decision Making:
Egocentric Discounting and Reputation Formation.'' \emph{Organizational
Behavior and Human Decision Processes} 83(2): 260-281.
\url{https://doi.org/10.1006/obhd.2000.2909}.

}
\clearpage
\section{Appendix A. Beliefs, acquisition value, and all information
sets}\label{appendix-a.-beliefs-acquisition-value-and-all-information-sets}

\subsection{A.1 Probability
identities}\label{a.1-probability-identities}

The primitive joint mass is

\[P(T=t,\theta,s,z)=\frac{P(T=t)}{2}
 p_t^{1\{s=\theta\}}(1-p_t)^{1\{s\ne\theta\}}
 r^{1\{z=\theta\}}(1-r)^{1\{z\ne\theta\}}.\]

Conditional on contradiction, the correct-\(z\) mass is \((1-p)r\), the
correct-\(s\) mass is \(p(1-r)\), their sum is \(d\), and their
difference is \(h\). Thus the accuracy advantage of reversal is
\(g=h/d\). Conditional on confirmation, posterior accuracy is

\[p_F=\frac{pr}{pr+(1-p)(1-r)}>p>1/2.\]

The type posteriors \(m_C,m_F\) obey \(dm_C+(1-d)m_F=\mu\). Since
\(p_H>p_L\) and \(r>1/2\), \(d_H<d_L\) and \(m_C<\mu<m_F\).

\subsection{A.2 What sequential consistency pins
down}\label{a.2-what-sequential-consistency-pins-down}

Before acquisition, \(c\) is independent of \((T,\theta,s)\). Every
admissible acquisition strategy or perturbation is a function of
\((s,c)\) alone. For an acquired history, \(c\) conditional on \(s,e,z\)
is therefore independent of \(T\): its density is proportional to
\(f(c)P(e=1\mid s,c)\), with the type-dependent likelihood of \((s,z)\)
canceling from its normalization. Any subsequent action strategy depends
only on \((s,e,z,c)\) and independent mixing, so it has no incremental
type information conditional on the complete record. The audit flag,
conditional on action, is also type-independent by construction.
Consequently audited \(e=0\) histories have posterior \(\mu\) and
audited acquired histories have posterior \(m_C\) or \(m_F\), even in
limits approaching an otherwise off-path action.

This conclusion does not impose type-independent actions when the
manager actually knows type. It uses the manager's feasible information
in this model. It also does not claim uniqueness in the unrestricted set
of weak PBE assessments with arbitrary, structurally incompatible
zero-probability beliefs. Such assessments are excluded by consistency,
not selected according to favorable welfare.

\subsection{A.3 Initial, uninformed, confirming, and contradictory
decisions}\label{a.3-initial-uninformed-confirming-and-contradictory-decisions}

Every final observation is a coarsening of histories with primitive
posteriors \(\mu,m_C,m_F\). Bayes-consistent perturbations preserve
their convex hull, including at off-path observations. The maximum
compatible difference in reputation payoff is consequently
\(\lambda B\). Condition (D) thus strictly rules out reversal without
evidence. The confirmation accuracy advantage exceeds \(2p-1\), so
confirming reversal is also strictly dominated. This argument holds
before deriving any equilibrium posteriors and therefore excludes
perverse action pooling independently of an off-path belief selection.

If no cost type acquired, a sufficiently cheap manager could acquire and
implement the accuracy-improving signal. Unacquired reputation on path
is \(\mu\). Audited acquisition has mean posterior \(\mu\) by the
identity in A.1. Assigning the lowest compatible posterior \(m_C\) to
every unaudited reversal bounds the reputation loss by
\(L=\lambda d(\mu-m_C)\), yielding the profitable deviation used in
Theorem 1. Thus \(q>0\) at each diagnostic label.

Given positive acquisition and no dominated reversals, any reversal
identifies contradiction. If acquired contradictions were not corrected
with positive probability, there would be no benefit to acquisition. If
contradictory actions mixed, their continuation difference would be
zero. The conditional acquired-confirming action is the same action as
without acquisition, and complete records have mean ability posterior
\(\mu\). Integrating over the two evidence realizations then gives (6),
equal to \(-c\) at a contradictory indifference. Since \(F\) has no atom
at zero, positive acquisition requires strictly positive continuation
value and full correction. The cost draw is sunk and gives no
information about the state, so it cannot justify a different
contradictory action across cost types.

\subsection{A.4 Direct calculation of acquisition value under targeted
audit}\label{a.4-direct-calculation-of-acquisition-value-under-targeted-audit}

Given a hypothesized \(q\) and full use, unacquired reputation is

\[(1-\tau)m_N+\tau\mu.\]

Acquired reputation, averaged over the signal, is

\[dm_C+(1-d)[(1-\tau)m_N+\tau m_F].\]

Subtracting the former from the latter and using A.1 gives

\[-(1-\tau)d(m_N-m_C).\]

A reversal audit does not change \(m_C\) and hence does not enter the
calculation. Acquisition supplies gross accuracy gain \(h\). Subtracting
its cost gives cutoff (4), and the equilibrium contradiction gain is

\[g-(1-\tau)\lambda D(q)=c^*/d>0.\]

This verifies the correction deviation that remains after the
acquisition threshold is solved. The other action deviations were
excluded in A.3. At \(c=c^*\), either acquisition choice is optimal, but
this cost has zero probability and has no effect on the allocation or
audit optimum.

\subsection{A.5 Existence with a general cost
distribution}\label{a.5-existence-with-a-general-cost-distribution}

Let \(B(q)=F[h-(1-\tau)L/(1-dq)]\), with the CDF extended by zero below
zero and one above \(C\). The map is continuous and nonincreasing. Since
(D) implies \(L<h\), \(B(0)>0\). Also \(B(F(h))\leq F(h)\), with a
strict inequality when \(\tau<1\). Hence \(q-B(q)\) is strictly
increasing and crosses zero exactly once. At that fixed point the
posteriors and strict action best responses in A.4 construct an
equilibrium. The acquisition and audit comparative statics in (7) follow
by differentiation; existence and uniqueness do not require
differentiability at every cost.

\section{Appendix B. Closed forms and global policy
verification}\label{appendix-b.-closed-forms-and-global-policy-verification}

For uniform cost, the smaller root of the no-audit equation is

\[q_0=\frac{C+hd-\sqrt{(C-hd)^2+4CdL}}{2Cd}.\]

The other root lies outside the economically relevant interval. Since
\(L<h\), the smaller root is positive. The inverse policy (10) has
derivative

\[\tau'(q)=\frac{C(1-dq)+d(h-Cq)}{L}>0.\]

Thus implementing larger \(q\) requires strictly greater retention
auditing. Differentiating \(H=\tau(1-dq)\) produces (14) and then (12).
A real-root solution outside \([q_0,q_f]\) is infeasible and must not be
used. Under the regularity condition in Theorem 3, the derivative
polynomial is strictly decreasing on this feasible interval whenever the
policy is interior, so its relevant root is the larger algebraic root;
the smaller root is outside the interval or a local minimum outside the
concave feasible region.

At \(\kappa=0\) welfare increases up to \(q_f\) and only complete
retention auditing implements that endpoint. For every \(\kappa>0\) in
Theorem 3's region, \(H'(q_f)>0\) and \(V'(q_f)<0\), so a slight
reduction from complete auditing increases welfare. This proves that
partial audit is not inferred from a parameter sweep.

For arbitrary admissible parameters without \(x_0<x_c\), the cubic can
be nonconcave and may have more than one local extremum. Theorem 2
remains sufficient for global optimization. It does not select a
favorable equilibrium: the acquisition allocation is already unique
under (D). Policy ties in the cubic problem are different from
signaling-equilibrium multiplicity.

\section{Appendix C. Welfare accounting and excluded
instruments}\label{appendix-c.-welfare-accounting-and-excluded-instruments}

Under any equilibrium, iterated expectations gives \(\mathbb E[m]=\mu\).
With linear reputation payoff, mean reputation is unchanged by audit
policy even though ex post contradicted managers can prefer concealment.
The real investigation cost of a uniform cutoff is

\[\int_0^{Cq}c\,\frac{dc}{C}=Cq^2/2.\]

This proves the total-surplus formula (11). A policy that orders
additional useful evidence but charges its real investigation cost to
someone outside the firm does not make that cost disappear from welfare.
Conversely, a profit calculation that omits it must not be labeled the
social optimum.

If the planner could directly buy the corrective evidence for every cost
type below \(h\), reveal it without an audit infrastructure cost, and
enforce its use, the feasible set would be different and could attain
\(q_f\) at lower cost. The article does not rule out that intervention.
The solved problem concerns environments where the relevant diagnostic
work is privately chosen and verification of its trail consumes
resources. Similarly, allowing sanctions, nonlinear contracts,
endogenous auditor competence, selective deletion, or a direct sorting
payoff changes the design problem. None is silently optimized by
Theorems 2 or 3.

\end{document}
